\section{Interpretation of Findings}
\label{sec:interpretation}

% ============================================================
\subsection{The Verdict on ``Defense Wins Championships''}
\label{sec:verdict}
% ============================================================

The central empirical question of this paper -- whether defensive dominance is
the primary pathway to championship success in the modern NFL -- receives an
unambiguous answer. All three pre-registered hypotheses are rejected, as
summarized in \Cref{tab:hypothesis-summary}. The data do not support the claim
that ``defense wins championships'' in any of its testable formulations: defense-dominant
teams do not win more playoff games than offense-dominant teams (H1 rejected),
Super Bowl winners are not drawn disproportionately from defense-dominant
archetypes (H2 rejected), and defensive \xvoa{} does not exhibit a stronger
association with playoff success than offensive \xvoa{} (H3 rejected). The
rejection is not marginal or dependent on a particular statistical procedure; it
is consistent across all five tests, both descriptive partitioning strategies,
and both era subsamples examined in \Cref{ch:results}. The popular narrative
that defensive excellence is the defining characteristic of championship teams
lacks empirical support in the 2018--2024 NFL.

% ============================================================
\subsection{What the Data Actually Show}
\label{sec:what-data-show}
% ============================================================

The evidence against the defensive narrative operates at every analytical level
employed in this paper. At the correlational level, the Spearman rank
correlation between offense share and playoff wins is positive and significant
($\rho = 0.246$, $p = 0.0167$, $n = 94$ playoff teams), directly contradicting
the directional prediction of the defensive hypothesis (\Cref{sec:spearman-results}).
The sign is unambiguous: teams whose quality is proportionally more
offense-driven win more postseason games, not fewer.

At the variance-explained level, the asymmetry is even more pronounced. Offensive
\xvoa{} per game accounts for 46.4\% of the variance in season win percentage,
compared to 14.2\% for defensive \xvoa{} -- a ratio of approximately $3.3 : 1$ in
favor of offense (\Cref{sec:pearson-results}). This finding means that if one
were forced to evaluate a team's competitive quality using only one side of the
ball, offense would provide more than three times as much information as defense.
The result is not an artifact of the \gds{} framework's construction; it reflects
the empirical reality that offensive execution varies more meaningfully across
teams and is more strongly coupled to winning than defensive execution.

At the effect-size level, the practical magnitude of the offensive advantage is
large. Cohen's $d = 0.672$ [0.574, 0.787] quantifies the standardized
difference in playoff wins between offense-dominant ($n = 143$) and
defense-dominant ($n = 29$) team-seasons across the full sample
(\Cref{sec:cohens-d-results}). In concrete terms, offense-dominant teams
averaged 0.60 playoff wins per season, compared to 0.00 for
defense-dominant teams -- no defense-dominant team-season in the study window
produced any playoff wins at all. This is not a subtle statistical
distinction but a difference visible to the naked eye of any attentive observer.

At the absolute level, the finding is starker still. Among the 29 defense-dominant
team-seasons in the study window (offense share $< -0.3$), zero won any
playoff games, and zero reached the Super Bowl. This null result stands against
94 total playoff team-seasons across seven seasons -- a sample large enough that
the complete absence of deep defensive postseason runs cannot reasonably be
attributed to chance. The logistic regression reinforces this pattern: both
offensive ($p = 0.003$) and defensive ($p = 0.027$) \xvoa{} are statistically
significant predictors of Super Bowl victory (\Cref{sec:logistic-results}), but
all seven Super Bowl champions during the study window had positive offense
share. The convergence of correlational, variance-explained, effect-size, and
absolute measures all pointing in the same direction constitutes a robust and
internally consistent empirical conclusion: offensive dominance is positively
associated with playoff success at every level of analysis.

% ============================================================
\subsection{The Nuanced Interpretation: Offense Gets You There, Defense Is the Final Filter}
\label{sec:nuanced-interpretation}
% ============================================================

The rejection of the ``defense wins championships'' hypothesis should not be
misinterpreted as a claim that defense is irrelevant to winning championships.
The data suggest a more refined structural interpretation: offensive quality is
the primary engine of playoff qualification and advancement, while competent
defense serves as a necessary -- but not sufficient -- condition at the championship
margin. The distinction between being the \emph{driver} of success and being a
\emph{filter} for the final stage is crucial. Offense determines which teams
reach the deep postseason; among teams that arrive there, defensive adequacy
helps determine which survives the final rounds.

Kansas City's 2023 Super Bowl victory illustrates this ``balanced elite''
archetype. Often characterized in popular discourse as a ``defensive''
championship, the \gds{} decomposition reveals a more nuanced picture:
the 2023 Chiefs had a positive offensive \xvoa{} ($+2.641$/game) alongside a
positive defensive contribution ($+1.896$/game), yielding an offense share of
$+0.581$ -- the lowest among Super Bowl winners but still firmly
offense-dominant. The popular narrative that Kansas City won ``because of their
defense'' is not supported by the expected-point framework; the offense
generated more absolute value than the defense. What distinguishes this
championship is that defensive quality contributed more \emph{relative to other
champions}, not that it was the primary engine of success.

The 2024 Philadelphia Eagles represent the opposite end of the champion
spectrum (offense $+7.740$/game, defense $+0.639$/game, share $= +0.923$),
demonstrating that championship teams need not have elite defense if their
offensive surplus is sufficiently dominant.

The concept of a ``competent defense threshold'' emerges naturally from the
quadrant analysis (\Cref{sec:quadrant-findings}). Among the top-quartile \gds{}
teams, the \textit{Offense Only} quadrant produced four Super Bowl winners
(19.0\% win rate) compared to two for \textit{Defense Only} (9.5\%), despite
identical sample sizes ($n = 21$). The implication is not that championship
teams need elite defense; it is that they need non-catastrophic defense. Five of seven champions recorded \emph{negative} defensive
\xvoa{}/game, meaning their defenses performed below league-average
expectations at the drive level. A team with a dominant offense and below-average
defense (Kansas City 2022, with defensive \xvoa{} of $-3.826$/game) can win the
Super Bowl. A team with elite defense and below-average offense has not done so
in seven seasons of data. The threshold is asymmetric: the minimum offensive
requirement for championship contention is substantially higher than the minimum
defensive requirement.

% ============================================================
\subsection{Why Offense Matters More: Structural Explanations}
\label{sec:structural-explanations}
% ============================================================

The empirical dominance of offense in predicting team success is not a statistical
curiosity but a reflection of structural features of modern professional American
football. Three mechanisms help explain why offensive quality is the stronger
differentiator.

First, offensive production exhibits higher between-team variance than defensive
production. The component decomposition in \Cref{sec:season-correlation} shows
that the standard deviation of offensive \xvoa{}/game across team-seasons exceeds
that of defensive \xvoa{}/game. Higher variance in a predictor variable
mechanically produces stronger correlations with outcomes, because the variable
provides more discriminating information between observations. In practical terms,
the gap between the best and worst offenses in the NFL is wider than the gap
between the best and worst defenses, which means that offensive quality is a
more informative signal for distinguishing good teams from poor ones. This
variance asymmetry is not a measurement artifact; it reflects the structural
reality that offensive schemes are more diverse, quarterback quality varies more
dramatically across rosters, and the passing game creates more opportunity for
separation between teams than the constrained action space of defensive play.

Second, the post-2018 NFL rule environment systematically amplifies offensive
variance. Increased enforcement of pass interference, roughing-the-passer
penalties, and restrictions on defensive contact beyond five yards \parencite{pereira2018roughing,barnwell2019passingrevolution} have expanded
the ceiling for offensive execution while compressing the floor for defensive
impact. Even the most elite defensive units in the modern game cannot prevent
completions to the degree that was possible before 2018, because the rules
constrain the tools available to them. The era analysis in \Cref{sec:era-findings}
shows a modest increase in the Pearson correlation between offensive \xvoa{} and
win percentage from 0.621 (2018--2020) to 0.739 (2021--2024), directionally
consistent with this structural explanation. The defensive Pearson coefficient
shows no corresponding increase (0.397 to 0.364), suggesting that defensive
execution has become less -- not more -- predictive of outcomes over time.

Third, offense possesses a structural ceiling advantage. The best offenses can
score on nearly any opponent in any game; the best defenses cannot prevent all
scoring against any opponent. This asymmetry creates a fundamental advantage
for offensive investment: a team that builds an elite offense creates a floor
of competitiveness that is higher than the floor created by building an elite
defense, because offensive points are under greater direct control than prevented
points. Defense is inherently reactive -- it responds to offensive formation,
personnel, and play-calling -- while offense is proactive, choosing the terms of
engagement. This proactive-reactive asymmetry means that schematic innovation
accrues disproportionately to the offensive side, where coordinators can design
novel concepts that defenses have not yet adapted to, while defensive innovation
is constrained to adjustments within the opponent's chosen framework.

% ============================================================
\subsection{Why the Myth Persists}
\label{sec:myth-persistence}
% ============================================================

If the empirical evidence so clearly favors offense, why does the belief that
``defense wins championships'' retain such cultural force in American football
discourse? Several cognitive and structural factors conspire to sustain the
narrative despite its empirical weakness.

The most powerful is selection bias in memorable outcomes. The Super Bowls that
dominate collective memory are disproportionately those with spectacular defensive
performances: the 2013 Seattle Seahawks dismantling Peyton Manning's record-setting
offense, the 2015 Denver Broncos winning with an ageing quarterback behind a
dominant defense. These games are memorable precisely because they are unusual -- a
defense-dominant team winning the championship is a departure from the base rate,
which makes it cognitively salient. The seven offense-dominant champions in this
study's window do not generate comparable narrative drama; a team with the best
offense winning the Super Bowl is unremarkable, and unremarkable events do not
anchor in memory. This is a textbook instance of availability bias \parencite{tversky1973availability}: the ease
with which defensive championship performances come to mind is mistaken for their
frequency.

The ``denominator problem'' compounds this selection bias. Observers readily
notice when a defense-dominant team wins a championship because the event is
visible and celebrated. They do not notice the far more frequent occurrence of
defense-dominant teams failing to reach the championship stage, because absences
are invisible. The data in \Cref{sec:archetype-findings} make this absence
concrete: only six defense-dominant teams (offense share $< 0$) even made the
playoffs across seven seasons, and none won more than a single game. These quiet
exits receive no media attention, no retrospective analysis, and no cultural
narrative. The myth persists in part because the denominator -- the set of
defense-dominant teams that failed -- is never examined with the same attention
as the rare numerator event of a defense-dominant champion.

Additionally, the narrative possesses a romantic appeal that transcends evidence.
The David-versus-Goliath framing -- a tough, physical defensive team outmuscling
a flashy, high-scoring offense -- resonates with cultural values of effort,
grit, and determination. It suggests that talent and scheme can be overcome by
intensity and will, a story more satisfying than the alternative: that the team
with better players executing a better offensive scheme typically wins. The
romantic narrative also benefits from confirmation bias in real-time game
viewing. When a team wins a low-scoring game, observers attribute the result to
defensive dominance rather than to the more prosaic explanations of opponent
offensive ineptitude, punting decisions, or game-script effects. Every defensive
stop is a narrative event; the absence of offensive production that made the stop
necessary goes unexamined.

The myth is not refuted by pointing to a single game or a single counter-example.
It is refuted by the systematic pattern across 224 team-seasons and seven
complete NFL cycles. The contribution of this paper is not to identify one
exception to the defensive narrative -- exceptions on both sides are easily
found -- but to demonstrate that the full weight of evidence, examined through
multiple statistical lenses and a novel measurement framework, consistently and
strongly favors offensive dominance as the primary pathway to championship
success in the modern NFL.
